\documentclass[11pt]{article}
\usepackage[T1]{fontenc}
\usepackage[utf8]{inputenc}
\usepackage[margin=1in]{geometry}
\usepackage{amsmath,amssymb,amsthm,booktabs,microtype}
\usepackage[hidelinks]{hyperref}

\newtheorem{theorem}{Theorem}[section]
\newtheorem{lemma}[theorem]{Lemma}
\newtheorem{proposition}[theorem]{Proposition}
\newtheorem{corollary}[theorem]{Corollary}

\title{Compatible orderings of binary vector spaces}
\author{Carptopus\\\small\texttt{carptopus@163.com}}
\date{Draft, 31 August 2026}

\begin{document}
\maketitle

\begin{abstract}
A total order on a finite vector space is called compatible when the unique
order-preserving bijection between any two subspaces of the same dimension is
linear. Cameron asked for the number of compatible orders on $V(n,q)$ and
observed that the answer was known only in dimensions one and two. We solve the
problem over the binary field. For every $n\ge4$, compatible total orders on
$\mathbb F_2^n$ form exactly eight orbits under $\mathrm{GL}(n,2)$, and
therefore their number is
\[
8|\mathrm{GL}(n,2)|=8\prod_{i=0}^{n-1}(2^n-2^i).
\]
Four orbits have zero as their least element. They are represented by two
explicit flag orders and by reversing the nonzero parts of those orders; the
remaining four are obtained by reversing the full order. The proof has one
finite computer-assisted base lemma: an exhaustive, symmetry-broken
classification in dimension four, checked by two implementations using
different invariants of ordered Fano planes. A local line condition then
forces an initial complete flag, and a uniqueness lemma shows that each of the
two relevant three-dimensional types has at most one lift in every higher
dimension. The counts in dimensions $1,2,3$ are respectively $2,12,10080$.
\end{abstract}

\section{Introduction}

Let $V=V(n,q)$ be an $n$-dimensional vector space over $\mathbb F_q$. A
total order $\prec$ on $V$ is \emph{compatible} if, whenever $U,W\le V$ have
the same dimension, the unique bijection from $U$ to $W$ that preserves their
induced orders is a linear isomorphism. Cameron introduced this condition
while studying monotone elements in inverse semigroups of partial linear
isomorphisms~\cite{cameron-notes}. Lexicographic orders give compatible orders
over every finite field. Cameron also determined the possible position of
zero, classified dimensions one and two, and noted that over $\mathbb F_2$
every order with zero extremal is compatible in dimension three. The counting
problem was restated as Problem~30.7 of the BCC30 problem collection, where the
answer was recorded as known only for $n=1,2$~\cite{bcc30}.

We give a complete answer for $q=2$. The binary case is rigid for two related
reasons. First, the one-dimensional condition forces zero to be the least or
greatest vector. Second, an ordered three-dimensional subspace is an ordered
Fano plane, so its linear type can be encoded by the positions of its seven
lines. An exhaustive four-dimensional base classification leaves four types
with zero least. Two of them satisfy a local condition saying that the first
three nonzero vectors in every three-space form a line; the other two are
obtained by reversing all nonzero vectors. The local condition forces the
initial segments of sizes $2^r-1$ to be subspaces. Once this flag is present,
the order outside its initial hyperplane is uniquely determined by the common
ordered Fano type.

The computer calculation is confined to dimension four. Searches in
dimensions five through seven were useful during discovery, but they are not
used to infer the general theorem. The general step is the local-to-global
flag argument and the unique-lift lemma proved below.

\section{Compatible orders and two explicit families}

Write vectors in coordinates $(x_1,\ldots,x_n)$ and identify a binary vector
with the integer $\sum_{i=1}^n2^{i-1}x_i$. For $x\ne0$, let
\[
h(x)=\max\{i:x_i=1\}.
\]
Define two total orders with zero first. In $L_n$, nonzero vectors are ordered
first by increasing $h(x)$ and then, within a fixed layer, by increasing key
\[
\lambda_h(x)=\sum_{i<h}2^{i-1}x_i.
\]
In $R_n$, the layers are again ordered by increasing $h(x)$, but
$\lambda_h(x)$ is decreasing inside each layer. Let $\rho$ keep zero first
and reverse the order of all nonzero vectors. Let $\iota$ reverse the full
total order.

For example, the nonzero parts in dimension four begin
\[
L_4=(1,2,3,4,5,6,7,8,9,10,11,12,13,14,15)
\]
and
\[
R_4=(1,3,2,7,6,5,4,15,14,13,12,11,10,9,8).
\]
After the standard $\mathrm{GL}(4,2)$ normalization by the first four
independent vectors, the second sequence becomes
\[
(1,2,3,4,5,7,6,8,9,11,10,14,15,13,12).
\]

\begin{proposition}\label{prop:explicit}
The orders $L_n,R_n,\rho L_n,\rho R_n$ and their full reversals are compatible
for every $n$.
\end{proposition}

\begin{proof}
Let $W\le\mathbb F_2^n$ have dimension $d$. Choose its reduced echelon basis
$w_1,\ldots,w_d$ using the highest nonzero coordinate as pivot, with the pivot
positions strictly increasing. The coordinate map
\[
\phi_W:\sum_{i=1}^dc_iw_i\longmapsto(c_1,\ldots,c_d)
\]
is linear. The highest nonzero coordinate of a linear combination is the
pivot belonging to the largest index $i$ with $c_i=1$. If two combinations
have the same largest active pivot, their highest differing lower coordinate
is the pivot belonging to the largest lower coefficient on which they differ.
It follows that $\phi_W$ carries the order induced by $L_n$ to $L_d$ and the
order induced by $R_n$ to $R_d$.

For two $d$-subspaces $U,W$, the order-preserving bijection is therefore
$\phi_W^{-1}\phi_U$, hence is linear. Reversing the nonzero parts of the
induced orders on both subspaces leaves the same rank-matching bijection, and
reversing both full orders does the same. Thus all eight displayed orders are
compatible.
\end{proof}

\section{The four-dimensional base classification}

The next lemma is the only computer-assisted component.

\begin{lemma}[Four-dimensional base]\label{lem:four}
With zero fixed as the least element, compatible total orders on
$\mathbb F_2^4$ form exactly four $\mathrm{GL}(4,2)$-orbits, represented by
\[
L_4,\quad R_4,\quad \rho L_4,\quad \rho R_4.
\]
\end{lemma}

\begin{proof}[Certificate and proof method]
An ordering of the seven nonzero points of $\mathbb F_2^3$ is classified up
to $\mathrm{GL}(3,2)$ by the set of position triples whose vectors sum to
zero. This line-position code is complete because a bijection of the Fano
plane that preserves its lines is linear. The $7!$ orders split into
$7!/168=30$ types.

For an order of the 15 nonzero vectors of $\mathbb F_2^4$, global linear
symmetry is removed by sending the first four independent vectors to
$1,2,4,8$. For each of the 30 possible ordered Fano types, a prefix search
adds vectors one at a time. Every hyperplane prefix must remain extendible to
that same Fano type. This is exhaustive after the stated symmetry breaking.
It gives the four normalized sequences
\begin{center}
\begin{tabular}{l}
\texttt{1 2 3 4 5 6 7 8 9 10 11 12 13 14 15}\\
\texttt{1 2 3 4 5 7 6 8 9 11 10 14 15 13 12}\\
\texttt{1 2 4 7 8 11 13 14 9 10 12 15 5 6 3}\\
\texttt{1 2 4 7 8 11 13 14 15 12 10 9 6 5 3}.
\end{tabular}
\end{center}

The first implementation compares hyperplane-coordinate orbit prefixes. A
second implementation shares the first-independent-basis normalization and
prefix-search architecture, but independently implements the finite-field
operations and compares Fano line-position prefixes directly. Both return
exactly these four sequences. This is a cross-check of the invariant
implementation, not a fully independent search design.

In the first two types, the first three nonzero vectors in every
three-dimensional subspace form a two-dimensional subspace. In the last two
types, the corresponding statement holds for the last three nonzero vectors.
The first two types induce distinct ordered Fano types.
\end{proof}

The finite lemma can also be viewed recursively: in dimension four,
compatibility is equivalent to requiring all hyperplanes to have one common
ordered Fano type. Every binary order in dimension three with zero extremal
is itself compatible, so no lower-dimensional condition is omitted.

\section{From a local line to a global flag}

\begin{lemma}[Initial flag]\label{lem:flag}
Let $n\ge3$, and let
\[
p_1\prec p_2\prec\cdots\prec p_{2^n-1}
\]
order the nonzero vectors of $V=\mathbb F_2^n$. Suppose that in every
three-dimensional subspace, the first three nonzero vectors form a
two-dimensional subspace. Then, for every $1\le r\le n$, the first $2^r-1$
nonzero vectors form an $r$-dimensional subspace.
\end{lemma}

\begin{proof}
The assertion is immediate for $r=1$. The vectors $p_1,p_2$ are distinct and
therefore independent. For any $x\notin\langle p_1,p_2\rangle$, the first two
vectors in $\langle p_1,p_2,x\rangle$ are $p_1,p_2$. Its third vector must
complete their two-dimensional subspace, so it is $p_1+p_2$. Hence the first
three global vectors form a two-space.

Assume the first $2^r-1$ vectors form an $r$-space $E_r$. Let $p$ be the next
vector and set $E_{r+1}=E_r+\langle p\rangle$. For $y\notin E_{r+1}$ and
$0\ne u\in E_r$, consider
\[
S=\langle u,p,y\rangle.
\]
Because $y\notin E_r+\langle p\rangle$, one has
$S\cap E_r=\langle u\rangle$. Thus $u$ is the first nonzero vector of $S$,
and $p$ is the second: every other point outside $E_r$ occurs no earlier than
the globally first such point $p$. The local hypothesis forces the third
point of $S$ to be $u+p$. Consequently $u+p\prec y$ for every nonzero
$u\in E_r$ and every $y\notin E_{r+1}$. Together with $p\prec y$, the whole
coset $p+E_r$ precedes every vector outside $E_{r+1}$. Hence the first
$2^{r+1}-1$ vectors are exactly $E_{r+1}\setminus\{0\}$. Induction completes
the proof.
\end{proof}

\section{Unique lifting and the classification}

\begin{lemma}[Unique lift]\label{lem:lift}
Fix an ordered Fano type $T$ whose first three points form a line. For every
$n\ge3$, there is at most one $\mathrm{GL}(n,2)$-orbit of orders on
$\mathbb F_2^n\setminus\{0\}$ satisfying both:
\begin{enumerate}
\item every three-dimensional subspace induces type $T$;
\item for every $r$, the first $2^r-1$ vectors form an $r$-subspace.
\end{enumerate}
\end{lemma}

\begin{proof}
We induct on $n$. Suppose two candidate orders have initial hyperplanes
$E,E'$. By the induction hypothesis, a linear isomorphism identifies their
induced orders on these hyperplanes. Let $p,p'$ be the first vectors outside
$E,E'$. Extend the hyperplane isomorphism linearly by sending $p$ to $p'$.
We may therefore assume that both candidates agree on $E$, have the same $p$,
and are defined on the same vector space.

Take two vectors $p+u,p+v$ in the outside coset. Choose a two-dimensional
subspace $K\le E$ containing $u$ and $v$; if $u,v$ span a smaller space,
extend it to such a $K$. Put $S=K\oplus\langle p\rangle$. Since every
nonzero vector of $E$ precedes $p$, the first three nonzero vectors of $S$
are precisely the three points of $K$, in the same order in both candidates,
and the fourth is $p$. Both restrictions have type $T$. Their
order-preserving linear isomorphism therefore fixes two basis vectors of $K$
and fixes $p$, hence is the identity on $S$. In particular, the relative
order of $p+u$ and $p+v$ is the same in both candidates. This holds for every
pair in the outside coset, so the two full orders agree.
\end{proof}

\begin{theorem}\label{thm:main}
For every $n\ge4$, compatible total orders on $\mathbb F_2^n$ form exactly
eight orbits under $\mathrm{GL}(n,2)$. Representatives are
\[
L_n,R_n,\rho L_n,\rho R_n,
\iota L_n,\iota R_n,\iota\rho L_n,\iota\rho R_n.
\]
Consequently, the number of labelled compatible total orders is
\[
\boxed{8|\mathrm{GL}(n,2)|
=8\prod_{i=0}^{n-1}(2^n-2^i)}.
\]
\end{theorem}

\begin{proof}
The one-dimensional subspaces show that zero has the same rank in every
two-point subspace. Hence zero is globally least or greatest. Reverse the
full order if necessary and assume zero is least.

All four-dimensional subspaces have one common ordered linear type. By
Lemma~\ref{lem:four} this type is one of
$L_4,R_4,\rho L_4,\rho R_4$. If it is one of the last two, apply $\rho$ to
the global order. Every three-dimensional subspace is contained in a
four-dimensional one, so in the resulting order its first three nonzero
vectors form a line. Lemma~\ref{lem:flag} supplies a complete initial flag.

The common ordered Fano type is now one of the two distinct types induced by
$L_4$ and $R_4$. Lemma~\ref{lem:lift} gives at most one orbit above each
type, while Proposition~\ref{prop:explicit} supplies $L_n$ and $R_n$.
Undoing $\rho$ gives exactly four orbits with zero least. Full reversal gives
four with zero greatest.

Finally, the action of $\mathrm{GL}(n,2)$ on total orders is free. A linear
map preserving a finite total order fixes each vector in its position and is
therefore the identity. Every orbit has size $|\mathrm{GL}(n,2)|$, proving
the formula.
\end{proof}

\begin{corollary}
The numbers in dimensions one, two and three are
\[
2,\qquad 2\cdot3!=12,\qquad 2\cdot7!=10080,
\]
respectively.
\end{corollary}

\begin{proof}
Zero must be extremal. For $n=1$ there are two orders. For $n=2$, any
ordering of the three nonzero vectors is compatible. The same is true for the
seven nonzero vectors when $n=3$: the only nontrivial proper subspaces are
Fano lines, and every bijection between two binary two-spaces that fixes zero
is linear.
\end{proof}

\section{Reproducibility and scope}

The four-dimensional certificate is supplied by \texttt{classify\_n4.py},
the hyperplane-coordinate prefix search, and
\texttt{verify\_n4\_independent.py}, the line-position-prefix verifier. The
corresponding JSON result files have SHA256 hashes
\begin{center}
\begin{tabular}{l}
\texttt{32b5a7c748270696372f57e5df890c963a4aef09cf583bd9c5e95cc2a642663a}\\
\texttt{cb682eb6e88009f7cda3b8c127f7fc323ac30c6a7910a0c0886e6ae328c3753f}.
\end{tabular}
\end{center}
The first result was made deterministic by excluding wall-clock timing from
the certificate. Finite recursive searches in dimensions five through seven
guided the discovery of the unique-lift lemma, but no claim from those
searches is used in the manuscript or included in this certificate.

The theorem answers the BCC30 counting problem over $\mathbb F_2$. It does
not classify compatible orders for $q>2$, where zero need not be extremal and
each projective point contains more than one nonzero vector. Our literature
check covered Cameron's original notes, the BCC30 restatement, and direct
searches using the definition and the expected counting formula. We found no
source that states the binary classification. This records the checked
boundary and is not a claim that unpublished or unindexed parallel work
cannot exist.

\section*{AI-assisted research disclosure}

Carptopus is the responsible author and takes responsibility for the claims,
proofs, source use, and released artifacts. OpenAI Codex was used for
AI-assisted literature organization, proof development and stress testing,
exact verification, adversarial auditing, and manuscript preparation.

\begin{thebibliography}{9}

\bibitem{cameron-notes}
Peter J. Cameron,
\emph{Combinatorics of inverse semigroups}, expository notes.
\url{https://webspace.maths.qmul.ac.uk/p.j.cameron/csgnotes/invsemi.pdf}.

\bibitem{bcc30}
Peter J. Cameron (ed.),
\emph{Problems from BCC30},
arXiv:2409.07216v1 (2024), Problem~30.7.
\url{https://arxiv.org/abs/2409.07216}.

\end{thebibliography}

\end{document}
